\documentclass[twoside]{article}
\usepackage{amsmath}
\usepackage{lipsum} % Package to generate dummy text throughout this template
\usepackage[english]{isodate}% http://ctan.org/pkg/isodate
\usepackage{pgfgantt}
\usepackage[sc]{mathpazo} % Use the Palatino font
\usepackage[T1]{fontenc} % Use 8-bit encoding that has 256 glyphs
\linespread{1.5} % Line spacing - Palatino needs more space between lines
\usepackage{microtype} % Slightly tweak font spacing for aesthetics
\usepackage{minted}
\usemintedstyle{autumn}
\usepackage[hmarginratio=1:1,top=32mm,left=25mm,right=25mm,columnsep=18pt,bottom=27mm]{geometry} % Document margins
\usepackage{multicol} % Used for the two-column layout of the document
\usepackage[hang, small,labelfont=bf,up,textfont=it,up]{caption} % Custom captions under/above floats in tables or figures
\usepackage{booktabs} % Horizontal rules in tables
\usepackage{float} % Required for tables and figures in the multi-column environment - they need to be placed in specific locations with the [H] (e.g. \begin{table}[H])
\usepackage{hyperref} % For hyperlinks in the PDF

\usepackage{lettrine} % The lettrine is the first enlarged letter at the beginning of the text
\usepackage{paralist} % Used for the compactitem environment which makes bullet points with less space between them
\usepackage{etoolbox}
\patchcmd{\thebibliography}{\section*{\refname}}{}{}{}
\usepackage{abstract} % Allows abstract customization
\renewcommand{\abstractnamefont}{\normalfont\bfseries} % Set the "Abstract" text to bold
\renewcommand{\abstracttextfont}{\normalfont\small\itshape} % Set the abstract itself to small italic text

\usepackage{titlesec} % Allows customization of titles
\renewcommand\thesection{\Roman{section}} % Roman numerals for the sections
\renewcommand\thesubsection{\Roman{subsection}} % Roman numerals for subsections
\titleformat{\section}[block]{\large\scshape\centering}{\thesection.}{1em}{} % Change the look of the section titles
\titleformat{\subsection}[block]{\large}{\thesubsection.}{1em}{} % Change the look of the section titles

\usepackage{fancyhdr} % Headers and footers
\pagestyle{fancy} % All pages have headers and footers
\fancyhead{} % Blank out the default header
\fancyfoot{} % Blank out the default footer
\fancyhead[C]{Naomi Pentrel $\bullet$ \href{mailto:n.pentrel@jacobs-university.de}{n.pentrel@jacobs-university.de} $\bullet$ \today} % Custom header text
\fancyfoot[RO,LE]{\thepage} % Custom footer text
\usepackage{graphicx}
\usepackage{wrapfig}

\usepackage{tocloft}

%\usepackage{csquotes}
%\MakeOuterQuote{"}

\hypersetup{
    bookmarks=true,
    unicode=false,
    pdftoolbar=true,
    pdfmenubar=true,
    pdffitwindow=false,
    pdfstartview={FitH},
    pdftitle={My title},
    pdfauthor={Author},
    pdfsubject={Subject},
    pdfcreator={Creator},
    pdfproducer={Producer},
    pdfkeywords={keyword1} {key2} {key3},
    pdfnewwindow=true,
    colorlinks=false,
    linkcolor=red,
    citecolor=green,
    filecolor=magenta,
    urlcolor=cyan
}

\usepackage{caption}
\usepackage{subcaption}
\usepackage[section]{placeins}
\usepackage{qtree}
\usepackage{gensymb}

\usepackage{setspace}

\usepackage{color}

%\doublespacing
% or:
%\onehalfspacing

%\renewcommand\cftchapafterpnum{\vskip10pt}
%\renewcommand\cftsecafterpnum{\vskip15pt}

%----------------------------------------------------------------------------------------
%	TITLE SECTION
%----------------------------------------------------------------------------------------

\title{\vspace{-15mm}\fontsize{24pt}{10pt}\selectfont\textbf{// All comments are NOT created equal}} % Article title


%----------------------------------------------------------------------------------------

\begin{document}
\thispagestyle{empty}
\pagenumbering{roman}
\begin{flushright}
    \includegraphics[scale=0.7]{Logo}
  \end{flushright}
  \vspace{20mm}
  \begin{center}
    \huge
    \textbf{Visualization of Theory Graphs of Semantically Annotated Documents in the Form of Spatial Presentations}
  \end{center}
  \vspace*{4mm}
  \begin{center}
   \Large by
  \end{center}
  \vspace*{4mm}
  \begin{center}
    \Large
    \textbf{Naomi Pentrel}
  \end{center}
  \vspace*{20mm}
  \begin{center}
    \large
    Bachelor Thesis Proposal in Computer Science
  \end{center}
  \vfill
  \begin{flushright}
    \large
    \begin{tabular}{l}
      
      \hline
      Prof. Dr. Michael Kohlhase \\
      \\
    \end{tabular}
  \end{flushright}
  \vspace*{8mm}
  \begin{flushleft}
    \large
    Date of Submission: \today \\
    \rule{\textwidth}{1pt}
  \end{flushleft}
  \begin{center}
    \Large Jacobs University Bremen - School of Engineering and Science
  \end{center}

\newpage
\noindent
  With my signature, I certify that this thesis has been written by me
  using only the indicates resources and materials. Where I have
  presented data and results, the data and results are complete,
  genuine, and have been obtained by me unless otherwise acknowledged;
  where my results derive from computer programs, these computer
  programs have been written by me unless otherwise acknowledged. I
  further confirm that this thesis has not been submitted, either in
  part or as a whole, for any other academic degree at this or another
  institution.

  \vspace{20mm}

    \includegraphics[scale=0.2]{Signature}
 \hfill Bremen, \today
  
\newpage

\thispagestyle{fancy} % All pages have headers and footers

%----------------------------------------------------------------------------------------
%	ARTICLE CONTENTS
%----------------------------------------------------------------------------------------

 \section*{Abstract}
Visualization of knowledge is important to foster learning. To optimize the representation of knowledge and its interdependencies, we will focus on identifying ways to present knowledge without loosing context. This research aims at taking an existing annotated corpus and presenting its contents in a way that allows the audience to see the dependencies of the covered topics. This form of presenting the knowledge should make it easy to examine materials that form the basis for the topics the audience is interested in.\\

Consider the typical lecture one attends as a student. Towards the end of the lecture students might have difficulties remembering earlier topics. When introducing a new topic it would be ideal to have a simple way to find the knowledge dependencies and present students with a simple way to catch up.\\

The OMDoc (Open Mathematical Documents) format is a content markup scheme for mathematical documents.\cite{OMDoc book} Using OMDoc, this research aims to visualize connecting some of what is written on Mathematics or other fields in a way that allows students to go back to the concepts they have forgotten to understand the topics currently covered in class. As such it will foster learning and could transform how we interact with course material as a whole. It is easy to see that this research would also have many other real world applications anywhere where knowledge is supposed to be transferred.\\ 

  \newpage
  \tableofcontents

  \clearpage
  \pagenumbering{arabic}

  \section{Introduction}




"Knowledge is the information necessary to support intelligent reasoning." \cite{Kohlhase}

In this paper we want to go form merely displaying information to representing knowledge by displaying a network/context the information is in. In this we will not try to semantically represent the meaning of the information that is shown but rather on showing the network, i.e. context the information is in.\cite{Probst} \\

-addresses fellow experts from your field (but not from your particular topic of research). 
- Here you should technically connect to the main concepts from that field and give an outline of your project, stating the research/engineering question that you want to get answered by your project.
  (target size: 1 page)
\newpage
  \section{Statement and Motivation of Research}
  This part should make clear which question, exactly, you are
  pursuing, and why your project is relevant/interesting. This is the
  place to cite relevant literature. Where does your project extend
  the state of the art? what weaknesses in known approaches to you
  hope to overcome? If you have carried out preliminary experiments,
  describe them here.

(target size: 3-4 pages)

\subsection{OMDoc}
(half a page)
\vspace*{8cm}
\subsection{Spatial Narration}
What is Prezi?
(half a page)
\vspace*{8cm}
\subsection{Combining OMDoc and Spatial Narration}
What is the idea?
(half a page)
\newpage
\subsection{Positive Effects of the Combination of Spatial Narration and OMDoc}
Why is this important for the world?
How will this improve learning?
(1 page)
\newpage
\subsection{Different Representations of Data}
What's the best representation one could do?
(1 page)
  (target size: 3-5 pages)
\newpage
\section{Planned Investigation}

  This is the technical core of the proposal. Here you lay out your
  plans of how you want to answer your research question specify your
  design of experiments or simulations, point out difficulties that
  you expect to encounter, etc.

How will this be implemented?
Using which frameworks?
  

\subsection{Forming Visualizations out of Knowledge Graphs which Convey Meaning}
(1 page)
\newpage

\subsection{Ordered Knowledge Graphs}
(1 page)
\newpage

\subsection{Meaningful Representation of Knowledge Graphs}
(1 page)
\newpage


\section{Evaluation Criteria}

????

  This section discusses criteria which can be used to evaluate the
  research results and which can be used to related research results
  to the already known state-of-the-art.

  (target size: 1-2 page)
\newpage
  \section{Timeline}

\textbf{Milestones}
\begin{itemize}
\item \textbf{Until February 15th}: Complete preliminary development to create a presentation from an XML
\item \textbf{Until April 1st}: Complete development of organized presentations that are generated from ordered knowledge graphs.
\item \textbf{Until May 1st}: Testing and Performance Measurement.
\item \textbf{Until May 10th}: Complete development of organized presentations that are generated from ordered knowledge graphs.
\end{itemize}


\begin{figure}[h]
\begin{center}
\begin{ganttchart}[y unit title=0.4cm,
y unit chart=0.5cm,
vgrid,hgrid, 
title label anchor/.style={below=-1.6ex},
title left shift=.05,
title right shift=-.05,
title height=1,
bar/.style={fill=gray!50},
incomplete/.style={fill=white},
progress label text={},
bar height=0.7,
group right shift=0,
group top shift=.6,
group height=.3,
group peaks={}{}{.2}]{24}
%labels
\gantttitle{Proposed Timeline}{24} \\
\gantttitle{December}{4} 
\gantttitle{January}{4} 
\gantttitle{February}{4} 
\gantttitle{March}{4} 
\gantttitle{April}{4} 
\gantttitle{May}{4} \\
%tasks
\ganttbar{Preliminary development}{2}{10} \\
\ganttbar{Graph visualization}{11}{16} \\
\ganttbar{Testing \& Performance}{17}{20} \\
\ganttbar{Final Report submission}{21}{21.5}

%relations 
\ganttlink{elem0}{elem1} 
\ganttlink{elem1}{elem2} 
\ganttlink{elem2}{elem3} 
\end{ganttchart}
\end{center}
\caption{Gantt Chart}
\end{figure}

  \section{Conclusions}

  Summarize the main aspects of the proposal.

  (target size: 1/2 page)

\newpage
\section{References}
\begin{thebibliography}{} % Bibliography - this is intentionally simple in this template
\bibitem {Kohlhase}
Kohlhase, Michael. "320441 CompLog Lecture Notes." (n.d.): n. pag. 17 Nov. 2014. Web. 19 Nov. 2014.
\bibitem {OMDoc book}
Kohlhase, Michael, and Alan Bundy. OMDoc - an Open Markup Format for Mathematical Documents. Berlin: SpringerLink, 2006. Print.
\bibitem {Probst}
Probst, Gilbert, Steffen Raub, and Kai Romhardt. \newblock Wissen Managen. \newblock Wiesbaden: S.n., 1997. Print.
\end{thebibliography}

\end{document}